
\section{Contraste bibliográfico: esfuerzo viscoso no equivale a fuerza viscosa}
\label{sec:ref26-ns}

La conservación de la masa de un medio continuo se expresa como \(\partial_t\rho+\nabla\cdot(\rho\mathbf{u})=0\). La condición \(\nabla\cdot\mathbf{u}=0\) se deduce cuando la densidad de cada partícula permanece constante, es decir, \(D\rho/Dt=0\) y \(\rho>0\); no debe introducirse sin justificarla para un gas compresible. Para densidad y viscosidad constantes y fluido newtoniano incompresible, sin otras fuerzas volumétricas,
\begin{equation}
 \nabla\cdot\mathbf{u}=0,\qquad
 \rho\,\frac{D\mathbf{u}}{Dt}=-\nabla p+\mu\nabla^2\mathbf{u}.
 \label{eq:ref26-ns-form}
\end{equation}
El término \(\mu\nabla^2\mathbf{u}\) es la \emph{divergencia del tensor viscoso}, no el tensor de esfuerzos viscosos \(\boldsymbol{\tau}=2\mu\mathbf{E}\) en sí. Una distribución de esfuerzos puede ser distinta de cero y poseer divergencia nula. Para viscosidad espacialmente variable se debe conservar \(\nabla\cdot(2\mu\mathbf{E})\), y no reemplazarlo sin más por \(\mu\nabla^2\mathbf{u}\).

\begin{tcolorbox}[title={Ejemplo guiado original: deformación plana}, studiumdef, breakable]
Considérese, lejos de fronteras sólidas, el campo bidimensional estacionario \(\mathbf{u}=(ax,-ay,0)\), con \(a\) constante. Su divergencia es \(a-a=0\), la aceleración material es \((a^2x,a^2y,0)\), y \(\nabla^2\mathbf{u}=\mathbf{0}\). La presión
\[
 p(x,y)=p_0-\frac{\rho a^2}{2}(x^2+y^2)
\]
satisface exactamente la ecuación de cantidad de movimiento. Pero
\[
 \mathbf{E}=\operatorname{diag}(a,-a,0),
 \qquad\boldsymbol{\tau}=2\mu\operatorname{diag}(a,-a,0)\ne\mathbf{0}.
\]
Por ejemplo, si \(a=2\,\mathrm{s^{-1}}\), \(\rho=1000\,\mathrm{kg\,m^{-3}}\) y \(\mu=0.001\,\mathrm{Pa\,s}\), los esfuerzos normales viscosos son \(+0.004\) y \(-0.004\,\mathrm{Pa}\), y en \((x,y)=(0.20,0.10)\,\mathrm{m}\) la presión es \(100\,\mathrm{Pa}\) menor que \(p_0\). Esto es una solución local idealizada, no una predicción sobre cómo se mantiene el campo en un recipiente finito.
\end{tcolorbox}

\textbf{Autocomprobación.} Calcule el rotacional del campo y distinga «flujo irrotacional» de «sin deformación viscosa». Discuta qué término reaparecería si \(\mu=\mu(x,y)\). Ejemplo y verificación algebraica propios, contrastados con los capítulos 4 de White \cite{ref26-white} y 11 de Shaughnessy, Katz y Schaffer \cite{ref26-shaughnessy}. Para la notación de gradiente y divergencia, véase Boas, capítulo 6 \cite{ref26-boas}.
